描述高能强子散射物理（例如在大型强子对撞机 LHC 上产生希格斯玻色子）的最大简化之一，是 R.~Feynman 引入的部分子模型~\cite{Feynman:1969ej}。按照这一图像，一个快速运动的强子（例如质子）可以被看作一束无相互作用的夸克和胶子（部分子），由其动量密度 $q(x)$ 和 $g(x)$ 表征，其中 $x$ 是部分子携带的纵向动量份额，$x= k^z/P^z$，而强子动量 $P^z \rightarrow \infty$。于是，涉及强子的硬散射截面可以表示为基本的部分子散射截面 $\hat \sigma$ 与部分子密度的卷积。在强相互作用的基本理论——量子色动力学（QCD）中，这一简单图像可以由所谓的因子化定理来论证~\cite{Mueller:1989hs,Sterman:1994ce,Collins:2011zzd}。量子场论引入的唯一修正是部分子密度依赖于方案和标度，后者可以通过重正化群方程加以研究~\cite{Altarelli:1977zs}。当然，部分子密度的方案和标度依赖会被部分子散射截面的类似依赖所抵消，从而使物理量在部分子的微扰定义下保持不变。

得益于渐近自由，部分子散射截面可以在 QCD 微扰论中计算，但部分子密度本质上是非微扰的。如上所述，Feynman 对部分子密度的定义是在无限动量系（IMF）中给出的，其中母强子具有无限大的动量~\cite{Feynman:1969ej}。这看起来是一个难以直观理解的数学极限。无限动量极限的确切含义实际上可以通过对微扰论中的 Feynman 图作 boost 来理解，正如 Weinberg 所做的那样~\cite{Weinberg:1966jm}。然而多年来，人们发现用光锥关联函数的形式表述部分子密度更为方便~\cite{Drell:1969ca}：引入光锥坐标，
\begin{equation}
      \xi^\pm = \frac{1}{2}(\xi^0\pm \xi^3) \ ,
\end{equation}
其他四维矢量同理，其中 $\xi^\mu$ $(\mu=0,1,2,3)$ 是时空坐标（把 $x$ 留给动量份额），强子沿 3 即 $z$ 方向运动。此时部分子密度计算为非定域关联符的矩阵元~\cite{Collins:1981uw}，
\begin{eqnarray}
   q(x, \mu^2) &=& \int \frac{d\xi^-}{4\pi} e^{-ix\xi^-P^+}  \langle P|\overline{\psi}(\xi^-)
   \gamma^+ \\
   && \times \exp\left(-ig\int^{\xi^-}_0 d\eta^- A^+(\eta^-) \right)\psi(0) |P\rangle \nonumber \ ,
\end{eqnarray}
其中 $\psi$ 是夸克狄拉克场，$A^\mu$ 是胶子势，$g$ 是强耦合常数，$\mu^2$ 是重正化标度。这里所有无限 boost 因子都消失了，IMF 物理完全通过被 boost 的夸克和胶子算符体现。上述矩阵元现在与强子外部动量 $P^\mu$ 无关！然而付出的代价是，部分子密度不再代表一个等时关联，而这正是 Feynman 当初在 IMF 中的原始定义。

尽管摆脱了无限 boost 已是一大进步，光锥关联函数绝非容易计算。例如，要用本质上为欧氏的 Wilson 格点 QCD 方法~\cite{Wilson:1974sk}，就必须去掉实时间依赖。人们必须对夸克场之间的间隔作 Taylor 展开，得到带若干阶导数的定域算符，其矩阵元与部分子分布的矩相联系~\cite{negele}。然而，测量含许多导数的算符矩阵元本质上噪声很大，因此只能计算前几个矩。多年来人们发明了多种克服这一困难的方法~\cite{Aglietti:1998mz,Liu:1999ak,Detmold:2005gg}，但没有一种对现实的数值模拟而言足够有前景。

另一种被一些人大力倡导的途径是光前量子化~\cite{Ji:1992ku,Burkardt:1995ct,Brodsky:1997de}，它与 IMF 中不含时微扰论的形式相关联~\cite{Weinberg:1966jm}。事实上，动力学理论的光前表述可以追溯到半个多世纪前 Dirac 的原始论文~\cite{Dirac:1949cp}。该框架大致如下：把 $\xi^+$ 取作新的时间，引入等``时''对易子来量子化场论。$\xi^+=0$ 处的场具有以 Fock 粒子表出的正则平面波展开式。新的哈密顿量为 $\hat H_{\rm LC}= \hat P^-$，可以用它像在非相对论量子力学中那样发展不含时微扰论。然后可以通过求解非微扰的类薛定谔方程确定 $P^-$ 的谱和本征态，
\begin{equation}
      \hat P^-|P\rangle = (M^2/2P^+)|P\rangle \ ,
\end{equation}
其中本征值 $M^2$ 是强子质量，$P^+$ 是光主动量。在 QCD 中，最方便的是采用 $A^+=0$ 规范，此时 Fock 粒子就是部分子。强子态可以用 Fock 展开表示~\cite{Brodsky:1997de}，
\begin{equation}
             |P\rangle = \sum_{n\alpha, \lambda_i} \int \Pi_i \frac{dx_id^2k_{\perp i}}{\sqrt{2x_i}(2\pi)^3} \psi_{n\alpha}(x_i, k_{\perp i}, \lambda_i) \left|n\alpha
             : x_iP^+, k_{\perp,i},\lambda_i \right\rangle \ ,
\end{equation}
其中 $n$ 是部分子数目，$\lambda_i$ 是螺旋度标记。所有部分子都具有纵向份额 $x_i$ 和横向动量 $k_{\perp, i}$。指标 $\alpha$ 对给定部分子数和部分子螺旋度下所有可能的振幅 $\psi_{n\alpha}$ 求和。质子和介子 Fock 态的例子可参见文献 \cite{Ji:2002xn,Ji:2003yj}。一旦有了光前波函数，就可以计算任何感兴趣的部分子物理观测量。

尽管理论有许多优点——例如真空变得``平庸''、洛伦兹群多出一个成为运动学的生成元等等~\cite{Brodsky:1997de}——在光锥面上求解场论是出了名地困难。除了二维理论和按 Fock 粒子数的微扰展开之外，3+1 维的非微扰计算没有找到系统的近似方法。普通的渐进场论其静态性质可以在欧氏格点上模拟；与之不同，光前理论不存在这样的形式体系。这一困难可能与它们本质上是闵氏型的事实有关。以横向格点表述光前理论的混合方案已被提出并探索过~\cite{Bardeen:1979xx}，但迄今尚未带来成功的模拟。

在最近的文献中，人们提出了一种利用欧氏格点 QCD 计算部分子分布的新方法~\cite{Ji:2013fga, Ji:2013dva, Xiong:2013bka,Hatta:2013gta,Lin:2014zya}。该方法本质上回到 Feynman 对部分子密度的原始定义，即从与强子内夸克场的空间关联函数相联系的普通动量分布 $n(\vec{k}, P^z)$ 出发。这个量可以用普通格点 QCD 计算，但它依赖于强子动量 $P^z$。Feynman 分布是在 $P^z\rightarrow \infty$ 极限下得到的，即
\begin{equation}
                q(x) \sim \lim_{P_z\rightarrow\infty} \int d^2k_\perp dk^z n(\vec{k},P^z)\delta (x-k^z/P^z) \ .
\end{equation}
然而，格点模拟无法直接给出无穷大 $P^z$ 的结果。那么如何从格点上可能实现的较大 $P^z$ 结果恢复 $P^z=\infty$ 极限呢？另一个微妙之处是场论通常有紫外（UV）发散。在一个典型的 Feynman 图中，UV 发散正规化的先后次序与 $P^z\rightarrow \infty$ 的次序彼此不可交换。定义部分子分布要求 $P^z$ 远大于截断标度，而格点模拟只有在 UV 截断远大于强子动量时才有意义。

文献 \cite{Ji:2013dva} 讨论了上述问题的解决办法：可在格点 QCD 中计算的 $n(\vec{k}, P^z)$ 可以通过微扰可算的乘法因子 $Z$ 加上 $(1/P^z)^2$ 幂次的高阶修正与部分子密度 $q(x)$ 相匹配。该方法可以推广到任何其他光前量，并可表述为存在一个（或多个）大动量标度 $P^z$ 时格点 QCD 的有效场论。这种情形与如今已被充分理解并获得广泛应用的重夸克有效场论（HQET）非常相似~\cite{Manohar:2000dt}。这里我较为详细地讨论这一新的{\it 大动量有效场论}（LaMET），它允许在有限 $P^z$ 下以系统近似在格点 QCD 中计算部分子物理。

让我们首先回顾 HQET 如何工作。考虑含一个重夸克的物理系统，例如价夸克组分为 $b\bar u$ 的 $B^-$ 介子。该系统的一个一般物理观测量 $O$（例如弱衰变常数）将依赖于 $b$ 夸克质量参数 $m_b$。由于渐近自由，对于大的 $m_b>\!>\Lambda_{\rm QCD}$（$\Lambda_{\rm QCD}$ 为 QCD 标度），$O(m_b)$ 可以按 $1/m_b$ 的幂次展开，
\begin{equation}
             O(m_b/\Lambda) = Z(m_b/\Lambda, \Lambda/\mu)o(\mu) + {\cal O}(1/m_b)+ ...
\end{equation}
这里我们假定 $O$ 依赖某个 UV 截断 $\Lambda$；在完整理论中经过恰当的重正化之后，它也可以换成重正化标度。展开的首项包含匹配系数 $Z$ 中重夸克质量的对数依赖，后者有关于 $\alpha_s$ 的微扰展开。$o(\mu)$ 是在有效理论（即 HQET）中定义的量，其中 $b$ 夸克实际上是无限重的。$\mu$ 是 HQET 中的重正化标度。次领头项至少被 $1/m_b$ 的一个幂次压低。

从微扰论很容易理解 HQET 在做什么。考虑外部波函数和内部传播子中含有夸克质量依赖的 Feynman 图。在完整理论中，这些图在恰当的 UV 正规化之后是完全确定的。在有效理论中，人们简单地在积分内部取 $m_b\rightarrow \infty$。这个极限不会改变所计算量的红外性质。因此，完整理论和有效理论中的 $O$ 与 $o$ 分别包含相同的非微扰物理。但它们的 UV 物理不同。这一差异体现在乘法匹配因子 $Z$ 上：它是红外自由的，因而可以在微扰论中计算。这里的非平凡之处在于 $Z$ 是乘法的。这可以像证明微扰可重正性那样逐阶在微扰论中证明~\cite{Collins:1984xc}。当然，我们忽略了可能存在几个相同量子数算符混合的复杂性，也忽略了在一定参数空间中 $Z$ 与 $o(\mu)$ 之间可能出现卷积的可能性。

HQET 应用中的一个重要物理观察是：一旦 $m_b$ 进入微扰区域，$O(m_b, \Lambda)$ 与其有效理论对应量 $o(\mu)$ 之间的差别就在掌控之中，也就是说我们知道如何计算它们之间的匹配。在某些情形下需要计入幂次修正，这些修正可以被分类并仔细研究。因此在一定程度上，重夸克的物理类似于无限重夸克的物理。什么时候可以把一个夸克视为重的？对某些物理观测量而言，质量为 1.5 GeV 的粲夸克已经是好的候选者；对质量为 4.2 GeV 的底夸克，预期 HQET 会工作得相当好。
现在可以考虑 LaMET 的构造。考虑一个依赖于大强子动量 $P^z$ 的格点（欧氏）观测量 $F$。利用渐近自由，我们可以系统地展开 $P_z$ 依赖，
\begin{equation}
             F(P^z/\Lambda) = Z(P^z/\Lambda, \Lambda/\mu)f(\mu) + {\cal O}(1/(P^z)^2)+ ... \ .
\label{Eq:Expansion}
\end{equation}
量 $f(\mu)$ 定义在 $P^z\rightarrow \infty$ 的理论中，正如 Feynman 部分子模型中那样。换言之，$f(\mu)$ 是一个光前关联或部分子观测量，包含所有的红外共线奇异性。因此，Feynman 的部分子模型事实上就是针对以大动量运动的核子的一种有效理论！这一展开的重要讯息是它可能在适度大的 $P^z$ 就已收敛，从而可以触及定义在无穷大 $P^z$ 处的量。通过精确计算匹配因子 $Z$ 及高阶修正，可以使部分子物理的提取更加精确。

上述展开的例子已在文献 \cite{Ji:2013dva, Ji:2013fga} 中讨论过。由于大动量标度来自外部态，我们不试图通过拉氏量框架来定义该有效理论，正如微扰 QCD 关于因子化定理的讨论那样。


格点准观测量的动量依赖可以通过重正化群来研究。通过
\begin{equation}
            \gamma(\alpha_s) = \frac{1}{Z}\frac{\partial Z}{\partial \ln P^z} \ .
\end{equation}
定义反常量纲。于是，
\begin{equation}
            \frac{\partial F(P^z)}{\partial \ln P^z} = \gamma(\alpha_s) F(P^z) + {\cal O}(1/(P^z)^2) \ ,
\end{equation}
准确到幂次修正。利用上式可以对涉及 $P^z$ 的大对数求和。


上述展开或有效描述得以存在的原因与 HQET 的存在类似。当在施加 UV 正规化之前先在 $F(P^z)$ 中取 $P^z\rightarrow \infty$ 时，按构造可以从 $\hat F$ 恢复光锥算符 $\hat f$。另一方面，格点矩阵元是在先施加 UV 正规化（格点截断）之后在大 $P^z$ 处计算的。于是矩阵元 $f$ 与 $F$ 之间的差别就是极限次序的问题。这正是有效场论的标准设置。不同的极限不改变红外物理。事实上，基于 Feynman 图的因子化可以像在重正化方案中那样逐阶证明。

上述有效场论展开给出了在格点 QCD 中研究部分子物理的一张处方：首先，从某个特定的部分子观测量 $\hat f$ 出发，它是由光锥场构成的算符。然后构造它的欧氏版本 $\hat F$，后者在无限洛伦兹 boost 下趋于 $\hat f$。在动量 $P^z$ 较大的强子中计算 $\hat F$ 的格点矩阵元，并用式~\ref{Eq:Expansion} 提取部分子物理 $\hat f$。当然，矩阵元 $F$ 既依赖于 $P^z$，也依赖于所有格点 UV 伪影。所有这些物理都将被吸收进匹配因子 $Z$。

前面的工作已经给出了 LaMET 如何作用于部分子密度的例子；这里我想讨论光前波函数本身。一旦知道完整的光前波函数，就可以计算任何部分子物理，而这正是光前量子化所追求的主要对象。为简单起见，考虑 $\pi^+$ 介子，尽管类似的讨论也适用于质子。领头的光前波函数与矩阵元
\begin{equation}
  \langle 0|\bar d(0)\gamma^+\gamma_5 u(\xi^-,\xi_\perp)|\pi^+(P)\rangle \ ,
\label{Eq:meson}
\end{equation}
相联系，其中 $u$ 和 $d$ 是坐标空间中的上、下夸克场，定义于 $A^+=0$ 规范。事实上，光前理论中所有的 Fock 波函数都可以表示成强子态与 QCD 真空之间的光前关联。上述算符的规范不变版本需要若干条沿 $\xi^-$ 方向通往 $\xi^- =+\infty$ 或 $-\infty$ 的规范链。$\xi^-$ 方向的选择必须一次且永远地固定下来；两种选择物理上不等价，彼此通过复共轭相联系~\cite{Ji:2002xn,Ji:2003yj}。为保证完全的规范不变性，无穷远处的规范链必须用横向规范链连接起来，因为在这一奇异规范中横向规范势在 $\xi^-=\pm \infty$ 处不为零。需要指明连接的方式，使上述矩阵元唯一确定。这可以用沿 $\vec{\xi}_\perp$ 的一条直规范链来实现~\cite{Ji:2002aa}。在多场关联的情形下，这些端点链可以有无限多种选取方式。在强子动量趋于无穷的某种极限过程中，这些端点链应当不作贡献，因为它们处在物理关联长度之外，下面将讨论这一点。

在 LaMET 中，我们需要找到一个准算符，它在无限洛伦兹 boost 下恢复式~\ref{Eq:meson} 中的算符。我们选取
\begin{equation}
      \hat F_\pm(z, \xi_\perp) = \bar d (0)W^\dagger(\pm\infty,0;0) \gamma^z\gamma_5 W(\pm\infty,z;\xi_\perp)u(z,\xi_\perp)
\end{equation}
其中所有场都位于 $\xi^0\equiv t = 0$ 并沿 $\xi^3\equiv z$ 方向分布，$\gamma^+$ 被 $\gamma^z$ 替换。为了获得规范不变性，我们对每个场插入一条沿 $z$ 方向延伸至 $z=\pm \infty$ 的规范链，
\begin{equation}
        W(\pm\infty, z; \xi_\perp) = P\exp\left(\mp ig\int^{\pm\infty}_z dz'A^z(z', \xi_\perp)\right) \ ,
\end{equation}
其中 $P$ 表示路径排序。在协变规范之类的规范中，$\hat F$ 已经是规范不变的。但在格点上，必须在较大的 $\pm z$ 处插入一条规范链以连接 $W^\dagger(\infty,0;0)$ 和 $W(\infty,z;\vec{\xi}_\perp)$。我们同样按光前量子化情形的需要定义末端规范链。再次强调，结果应当与这些端点链无关，因为物理关联长度是有限的。

在此，把从高能散射提取部分子物理的做法与 LaMET 作一比较是有益的。在高能散射中，人们从依赖某些大动量 $Q$ 的散射截面出发，进行标度分离并证明因子化定理，把这些截面与部分子性质联系起来；后者最终再从实验数据中提取~\cite{Gao:2013xoa}。另一方面，在欧氏格点上，人们计算依赖大强子动量 $P^z$ 的准观测量，可以用 LaMET 进行标度分离，把这些准观测量与部分子物理联系起来。这里有完全的类比。简言之，LaMET 使得用格点``数据''提取部分子物理成为可能。

正如同一部分子分布可以从不同的硬散射过程中提取出来，同一光前物理也可以从不同的格点算符中提取出来。给出相同光前物理的所有算符构成一个普适类。普适类的存在使人们可以探索不同的算符 $\hat F$，使得有限 $P^z$ 下的结果尽可能接近大 $P^z$ 下的结果。最近的一篇文献给出了探索不同算符以计算核子总胶子极化的例子~\cite{Hatta:2013gta}。

最后，评论一下本文方法与光前量子化的关系。用 Feynman 的原始语言，部分子物理由如下类型的矩阵元描述：
\begin{equation}
   \langle P_\infty|\hat F|P_\infty\rangle \ ,
\end{equation}
其中 $\hat F$ 是非定域、不含时的算符，$P_\infty$ 表示趋于无穷大的动量，可能带有残余动量，如同 Weinberg 论文中所用的那样~\cite{Weinberg:1966jm}。另一方面，可以写 $|P_\infty \rangle
= U_\infty|p_{\rm res} \rangle $，其中 $p_{\rm res}$ 是残余动量，而 $U_\infty$ 是希尔伯特空间中实现无限 boost 的幺正变换。于是可以把上面的量写成
\begin{equation}
   \langle p_{\rm res}|U^\dagger \hat FU|p_{\rm res}\rangle \ ,
\end{equation}
其中 boost 算符作用在 $F$ 上。这个无限 boost 将给出光前算符 $\hat f\equiv U^\dagger \hat FU$，
\begin{equation}
       \langle P_\infty|\hat F|P_\infty\rangle  = \langle p_{\rm res}|\hat f|p_{\rm res}\rangle \ ,
\end{equation}
其后一矩阵元不依赖残余动量。本文通过取大 $P^z$ 极限直接处理矩阵元 $\langle P_\infty|\hat F|P_\infty\rangle$，而光前量子化则试图用光前波函数计算该等式的右边。

上述讨论阐明了强子态时空图像的重要问题。人们常常画一个半径约 1fm 的圆来表示静止强子的坐标空间图像，但这在量子力学中会遇到问题，因为在平面波态中强子可以处于空间任何位置。可以构造波包来局域化强子，但在相对论理论中质心运动与内部运动之间便没有干净的分离，除非强子非常重。当考虑以一定速度运动的强子的空间图像时会出现更多问题，此时人们常画一个洛伦兹收缩的圆。强子态中长度收缩的物理意义往往是不清楚的。

\begin{figure}[!ht]
\begin{center}
\includegraphics[width=0.50\textwidth]{images/spacetime.pdf}
\end{center}
\caption{在高动量强子所在的参考系中沿空间方向 $z$ 的线段，变换为静止强子的光前坐标中沿光锥方向的线段，其长度增加了 boost 因子 $\gamma$；在 IMF 中 $\gamma$ 趋于无穷。}
\label{fig:cryosys}
\end{figure}


可以转而考虑平面波强子态中的场关联函数。一般地说，当强子静止（质心动量为零）时，场关联函数沿所有空间方向的关联长度约为 1fm $\sim \Lambda_{\rm QCD}$。超过这个范围，关联应逐渐降到零。现在考虑沿 $z$ 方向以速度 $v$ 运动的强子，引入洛伦兹因子 $\gamma=
1/\sqrt{1-v^2}$。价部分子将拥有固定的强子纵向动量份额 $x=k^z/P^z$。因此，价部分子的典型纵向动量将按 $xP^z$ 增长。相应地，纵向关联长度将表现为 $1/xP^z$，与 $\gamma$ 因子成反比。这可以被看作洛伦兹收缩效应。于是，要用 LaMET 计算价部分子分布，纵向空间分辨率必须随核子动量提高。如果只关心价部分子，纵向格点盒子尺寸也可以缩小 $\gamma$ 倍。这样格点总数便可与横向方向的相当。

另一方面，如果关心小 $x$ 部分子，就必须让强子动量足够大。能达到的最小 $x$ 为 $x\sim \Lambda_{\rm QCD}/P^z$。此时关联长度约为 $ 1/\Lambda_{\rm QCD}$，即强子的正常尺寸。例如在 $
x\sim 10^{-4}$ 处，强子动量必须在 3 TeV 左右，而普通价部分子的关联长度将为 $10^{-3}$fm。显然，要在如此小的 $x$ 处计算部分子密度，需要在 $z$ 方向数千个格点。

在 LaMET 中人们看到的是收缩的强子，而在光前计算中关联长度却急剧无界增长。这可以从图 1 看出，其中我们画出了沿 $z$ 轴的关联长度 $\ell$。在光前计算中，同样的关联被 boost 到沿 $\xi^-$ 方向，增大了 $\gamma$ 倍。因此对价部分子而言，光前坐标中的关联长度将是 $\sim 1/\Lambda_{\rm QCD}$ 的量级。对小 $x$ 部分子，关联长度可长达 $P^z/(M\Lambda_{QCD})$，其中 $M$ 是强子质量。因此，光前量子化计算无法限制在一个紧凑的时空区域内，这使得蒙特卡罗模拟变得困难。

总之，我们在本文中讨论了格点 QCD 模拟中大动量强子的新有效场论概念。该有效场论使人能够从强子动量为几个 GeV 量级的实际计算中提取光前部分子物理。相关展开包含微扰可算的匹配系数和高阶幂次修正。这使得格点数据在研究 QCD 束缚态性质方面与真实实验数据同样有用。

